\section{闭包与不可约性}\label{sec:closures-and-irreducibility}

本节的首要目标是讨论拟射影簇的不可约性和不可约分解. 有了对 Zariski 拓扑的新认识, 我们在拓扑框架内表述这些概念. 首先从拓扑空间不可约的一般含义说起.

\begin{Definition}{可约与不可约拓扑空间}\leavevmode\label{def:reducible-and-irreducible-topological-spaces}
若对某些真闭子集 $X_{1},X_{2}\subsetneq X$ 有 $X=X_{1}\cup X_{2}$, 则称拓扑空间 $X$ \emph{可约}. 若 $X$ 既非空又不可约化为上述形式, 则称它\emph{不可约}.
\end{Definition}

根据 Examples~\ref{ex:zariski-topology-on-projective-varieties} 和~\ref{ex:zariski-topology-on-affine-varieties} 对仿射簇和射影簇的 Zariski 闭集的描述, 若将 Definition~\ref{def:reducible-and-irreducible-topological-spaces} 应用于仿射簇和射影簇上的 Zariski 拓扑, 它便与我们先前关于仿射簇和射影簇不可约性的定义一致. 因此, Definition~\ref{def:reducible-and-irreducible-topological-spaces} 使我们能够把先前的概念推广到更一般的拟射影簇情形: 若拟射影簇关于其 Zariski 拓扑不可约, 则称它不可约.

注意, 不可约性对拓扑的选择颇为敏感. 例如, $\mathbb{R}=\mathbb{A}^{1}_{\mathbb{R}}$ 在 Zariski 拓扑下不可约, 因为不可能找到满足 $X_{1}\cup X_{2}=\mathbb{A}^{1}_{\mathbb{R}}$ 的仿射簇 $X_{1},X_{2}\subsetneq\mathbb{A}^{1}_{\mathbb{R}}$. 然而, $\mathbb{R}$ 在 Euclid 拓扑下可约, 因为 $(-\infty,0]$ 和 $[0,\infty)$ 是两个真闭子集, 且它们的并为 $\mathbb{R}$.

本节最后将证明, 每个拟射影簇都能唯一地写成不可约闭子簇的并. 这一结果的证明依赖于射影情形下不可约分解的存在性和唯一性 (Proposition/Definition~\ref{prop:projective-irreducible-decomposition}). 为了在拟射影情形和射影情形的不可约分解之间转换, 现在引入闭包和稠密性这两个一般的拓扑概念.

\begin{Definition}{闭包和稠密子集}\leavevmode\label{def:closure-and-dense-subset}
设 $Y$ 为拓扑空间, $X\subseteq Y$ 为子集. $X$ 在 $Y$ 中的\emph{闭包}记为 $\mathrm{cl}_{Y}(X)$, 它是 $Y$ 中所有包含 $X$ 的闭集的交. 若 $\mathrm{cl}_{Y}(X)=Y$, 则称 $X$ 在 $Y$ 中\emph{稠密}.
\end{Definition}

\BookPairedNoteLeft{
当环境拓扑空间 $Y$ 可由上下文明确时, 通常使用简写 $\overline{X}=\mathrm{cl}_{Y}(X)$.
}{
注意, $X$ 在 $Y$ 中的闭包是 $Y$ 的闭子集, 原因只是任意多个闭集的交是闭的. 这一闭包包含在 $Y$ 中每个包含 $X$ 的闭子集内.
}
换句话说, $\mathrm{cl}_{Y}(X)$ 是 $Y$ 中包含 $X$ 的最小闭子集. 展开定义还可发现, $X$ 在 $Y$ 中稠密, 当且仅当 $Y$ 的每个开子集至少包含 $X$ 的一个元素 (Exercise~\ref{ex:density-characterized-by-open-sets}). 换句话说, 若一个集合遍布 $Y$, 出现在每个邻域中, 则它在 $Y$ 中稠密. 例如, 实数的任意开区间无论多小, 都包含有理数, 由此推出 $\mathbb{Q}$ 关于 Euclid 拓扑在 $\mathbb{R}$ 中稠密.

现在来看 Zariski 拓扑中闭包和稠密性的例子.

\begin{Example}{$\mathbb{A}^{1}$ 中集合的闭包}\leavevmode\label{ex:closures-in-affine-line}
若 $X$ 是 $\mathbb{A}^{1}$ 中的有限集, 则它是仿射簇, 因而在 $\mathbb{A}^{1}$ 中闭, 所以 $\overline{X}=X$. 另一方面, 若 $X$ 无限, 则包含 $X$ 的闭集只有整个 $\mathbb{A}^{1}$. 因此, $\mathbb{A}^{1}$ 的每个无限子集都在 $\mathbb{A}^{1}$ 中稠密.
\end{Example}

\begin{Example}{射影闭包}\leavevmode\label{ex:projective-closure-of-arbitrary-subset}
若 $X\subseteq\mathbb{P}^{n}$ 是射影空间的任意子集, 则闭包 $\overline{X}\subseteq\mathbb{P}^{n}$ 称为 $X$ 的\emph{射影闭包}. $X\subseteq\mathbb{A}^{n}_{0}$ 是仿射簇的特殊情形, 已在第~\ref{sec:projective-closures}~节详加讨论. 对一般的子集 $X\subseteq\mathbb{P}^{n}$, 关于射影闭包有一个在下文中有用的关键事实: $\mathcal{I}_{\mathbb{P}}(X)=\mathcal{I}_{\mathbb{P}}(\overline{X})$, 下面说明理由.

包含关系 $\mathcal{I}_{\mathbb{P}}(\overline{X})\subseteq\mathcal{I}_{\mathbb{P}}(X)$ 成立, 只是因为在 $\overline{X}$ 上消没的任何多项式也必在较小的集合 $X$ 上消没. 为证明另一个包含关系, 只需注意若多项式 $f$ 不在 $\overline{X}$ 上消没, 它就不能在 $X$ 上消没. 否则 $\overline{X}\cap\mathcal{V}_{\mathbb{P}}(f)$ 是闭集, 且
\[
  X\subseteq\overline{X}\cap\mathcal{V}_{\mathbb{P}}(f)\subsetneq\overline{X},
\]
这与 $\overline{X}$ 为 $X$ 的闭包矛盾.

由 $\mathcal{I}_{\mathbb{P}}(X)=\mathcal{I}_{\mathbb{P}}(\overline{X})$ 可知, $X$ 的射影闭包可表示为
\[
  \overline{X}=\mathcal{V}_{\mathbb{P}}(\mathcal{I}_{\mathbb{P}}(X)).
\]
\end{Example}

\begin{Example}{仿射开片在射影空间中稠密}\leavevmode\label{ex:affine-patch-is-dense-in-projective-space}
考察仿射开片 $\mathbb{A}^{n}_{0}\subseteq\mathbb{P}^{n}$. 我们断言 $\mathbb{A}^{n}_{0}$ 的射影闭包为 $\mathbb{P}^{n}$. 为说明这一点, 首先注意
\[
  \mathbb{P}^{n}=\mathbb{A}^{n}_{0}\cup\mathcal{V}_{\mathbb{P}}(x_{0})
    =\overline{\mathbb{A}^{n}_{0}}\cup\mathcal{V}_{\mathbb{P}}(x_{0}).
\]
由于 $\mathbb{P}^{n}$ 不可约, 且 $\overline{\mathbb{A}^{n}_{0}}$ 与 $\mathcal{V}_{\mathbb{P}}(x_{0})$ 都是闭集, 它们中必有一个等于 $\mathbb{P}^{n}$. 已知 $\mathcal{V}_{\mathbb{P}}(x_{0})\ne\mathbb{P}^{n}$, 例如 $\mathcal{V}_{\mathbb{P}}(x_{0})$ 不包含 $[1:0:\cdots:0]$, 因而 $\overline{\mathbb{A}^{n}_{0}}=\mathbb{P}^{n}$, 所以 $\mathbb{A}^{n}_{0}$ 在 $\mathbb{P}^{n}$ 中稠密.
\end{Example}

把 Example~\ref{ex:affine-patch-is-dense-in-projective-space} 的论证推广到拓扑情形, 可见不可约拓扑空间的每个非空开集都稠密.

\begin{Proposition}{非空开集是稠密的}\leavevmode\label{prop:nonempty-open-in-irreducible-space-is-dense}
设 $X$ 为不可约拓扑空间. 若 $U\subseteq X$ 为非空开集, 则 $U$ 在 $X$ 中稠密.
\end{Proposition}
\begin{Proof}
设 $U$ 为不可约拓扑空间 $X$ 的一个非空开子集. 定义 $Z=X\setminus U$, 这是 $X$ 的一个闭子集. 注意 $X=\overline{U}\cup Z$. 由 $X$ 的不可约性, $X=\overline{U}$ 或 $X=Z$, 但 $X\ne Z$, 因为 $Z$ 是非空集合 $U\subseteq X$ 的补集. 因而 $X=\overline{U}$, 所以 $U$ 在 $X$ 中稠密.
\end{Proof}

拓扑情形下不可约性的一个有用性质是, 可以在稠密子集上判定不可约性, 下一个结果验证了这一点.

\begin{Proposition}{稠密性与不可约性}\leavevmode\label{prop:dense-subspace-irreducible-iff-ambient}
设拓扑空间 $Y$ 有稠密子集 $X\subseteq Y$. 则 $X$ 在子空间拓扑下不可约, 当且仅当 $Y$ 不可约.
\end{Proposition}
\begin{Proof}
设 $Y$ 为拓扑空间, 且 $X\subseteq Y$ 稠密. 我们证明双向叙述的两个方向.

假设 $X$ 不可约. 为了证明 $Y$ 不可约, 令 $Y=Y_{1}\cup Y_{2}$, 其中 $Y_{1},Y_{2}\subseteq Y$ 是闭集; 我们须证明 $Y_{1}=Y$ 或 $Y_{2}=Y$. 定义 $X_{1}=X\cap Y_{1}$ 和 $X_{2}=X\cap Y_{2}$. 注意 $X_{1},X_{2}$ 在 $X$ 的子空间拓扑中都闭, 且 $X=X_{1}\cup X_{2}$. 所以由 $X$ 的不可约性, $X=X_{1}$ 或 $X=X_{2}$. 这意味着 $X\subseteq Y_{1}$ 或 $X\subseteq Y_{2}$. 由于每个 $Y_{i}$ 在 $Y$ 中闭, 进一步得到 $\overline{X}\subseteq Y_{1}$ 或 $\overline{X}\subseteq Y_{2}$. 但 $X$ 在 $Y$ 中稠密, 故 $\overline{X}=Y$, 从而 $Y=Y_{1}$ 或 $Y=Y_{2}$, 如所需.

反过来, 假设 $Y$ 不可约. 为证明 $X$ 不可约, 令 $X=X_{1}\cup X_{2}$, 其中 $X_{1},X_{2}\subseteq X$ 是闭集; 我们须证明 $X_{1}=X$ 或 $X_{2}=X$. 由于取闭包与有限并可交换 (Exercise~\ref{ex:closure-commutes-with-finite-unions}), 有 $Y=\overline{X}=\overline{X_{1}}\cup\overline{X_{2}}$. 因为 $\overline{X_{1}}$ 和 $\overline{X_{2}}$ 在 $Y$ 中均闭, $Y$ 的不可约性意味着 $Y=\overline{X_{1}}$ 或 $Y=\overline{X_{2}}$; 不妨设 $Y=\overline{X_{1}}$. 由于 $X_{1}$ 在 $X$ 中闭, Exercise~\ref{ex:closure-in-subspace-is-restricted-ambient-closure} 蕴含 $X_{1}=\overline{X_{1}}\cap X$, 由此 $X_{1}=Y\cap X$, 即 $X_{1}=X$, 如所需.
\end{Proof}

最后, 给出本节的主要结果.

\begin{BookPropositionDefinition}{不可约分解}\leavevmode\label{prop:quasiprojective-irreducible-decomposition}
设 $X\subseteq\mathbb{P}^{n}$ 为非空拟射影簇. 则存在不可约闭子簇 $X_{1},\ldots,X_{r}\subseteq X$, 使得对任意 $i\ne j$ 有 $X_{i}\not\subseteq X_{j}$, 且
\begin{equation}\label{eq:quasiprojective-irreducible-components-union}
  X=\bigcup_{i=1}^{r}X_{i}.
\end{equation}
此外, 不可约闭子簇 $X_{1},\ldots,X_{r}$ 在不计次序时是唯一的. 我们称它们为 $X$ 的\emph{不可约分支}, 称 \eqref{eq:quasiprojective-irreducible-components-union} 为 $X$ 的\emph{不可约分解}.
\end{BookPropositionDefinition}
\begin{Proof}
设 $X\subseteq\mathbb{P}^{n}$ 为非空拟射影簇, 可写成 $X=Z\cap U$, 其中 $Z\subseteq\mathbb{P}^{n}$ 闭, $U\subseteq\mathbb{P}^{n}$ 开. 我们分别证明不可约分解的存在性和唯一性.

\textbf{(存在性)} 考察闭包 $\overline{X}\subseteq\mathbb{P}^{n}$. 由射影簇的不可约分解的存在性 (Proposition/Definition~\ref{prop:projective-irreducible-decomposition}), 存在不可约射影簇 $Z_{1},\ldots,Z_{r}\in\mathbb{P}^{n}$, 使得对任意 $i\ne j$ 有 $Z_{i}\not\subseteq Z_{j}$, 且
\[
  \overline{X}=\bigcup_{i=1}^{r}Z_{i}.
\]
若对某个 $i$ 有 $Z_{i}\cap X=\varnothing$, 则其他不可约分支的并将是 $\overline{X}$ 的一个包含 $X$ 的真闭子集, 矛盾. 因此对所有 $i$ 有 $Z_{i}\cap X\ne\varnothing$, 并由 $X_{i}=Z_{i}\cap X$ 定义 $X$ 的非空闭子簇. 由构造有
\[
  X=\bigcup_{i=1}^{r}X_{i},
\]
只需再证明每个 $X_{i}$ 都不可约, 且对 $i\ne j$ 有 $X_{i}\not\subseteq X_{j}$.

由于 $X=Z\cap U$ 且 $Z$ 闭, 有 $\overline{X}\subseteq Z$, 从而对每个 $i$ 有 $Z_{i}\subseteq Z$. 因 $X_{i}=Z_{i}\cap X=Z_{i}\cap(Z\cap U)$, 可知 $X_{i}=Z_{i}\cap U$. 由于 $U$ 开, 所以 $X_{i}$ 是不可约簇 $Z_{i}$ 的非空开子簇, 由 Proposition~\ref{prop:nonempty-open-in-irreducible-space-is-dense} 知它稠密. 由 Proposition~\ref{prop:dense-subspace-irreducible-iff-ambient}, $Z_{i}$ 的不可约性蕴含 $X_{i}$ 的不可约性. 再者, 若对某个 $i\ne j$ 有 $X_{i}\subseteq X_{j}$, 则取闭包得到 $Z_{i}\subseteq Z_{j}$, 这不可能, 因为它们是 $\overline{X}$ 的不同不可约分支. 因此, 拟射影簇的不可约分解存在.

\textbf{(唯一性)} 假设 $X_{1},\ldots,X_{r}$ 和 $X'_{1},\ldots,X'_{s}$ 是 $X$ 的两个不可约分解. 由于取闭包与有限并可交换 (Exercise~\ref{ex:closure-commutes-with-finite-unions}), 我们得到射影簇 $\overline{X}\subseteq\mathbb{P}^{n}$ 的两个分解:
\begin{equation}\label{eq:closures-of-two-quasiprojective-decompositions}
  \overline{X}=\bigcup_{i=1}^{r}\overline{X_{i}}
  =\bigcup_{i=1}^{s}\overline{X'_{i}},
\end{equation}
其中闭包均在 $\mathbb{P}^{n}$ 中取. 现在证明每个分支都不可约且没有一个包含在另一个中, 从而说明 \eqref{eq:closures-of-two-quasiprojective-decompositions} 中的两个分解均为不可约分解.

对每个 $i$, 注意 $\overline{X_{i}}$ 包含不可约簇 $X_{i}$ 作为稠密子集, 所以由 Proposition~\ref{prop:dense-subspace-irreducible-iff-ambient}, 每个 $\overline{X_{i}}$ 都不可约. 同理, 每个 $\overline{X'_{i}}$ 都不可约. 此外, 由于 $X_{i}$ 在 $X$ 中闭, 对每个 $i$ 有 $X_{i}=\overline{X_{i}}\cap X$ (Exercise~\ref{ex:closure-in-subspace-is-restricted-ambient-closure}), 因而不可能对某个 $i\ne j$ 有 $\overline{X_{i}}\subseteq\overline{X_{j}}$, 否则将推出 $X_{i}\subseteq X_{j}$, 与 $X_{1},\ldots,X_{r}$ 是 $X$ 的不可约分解的假设矛盾. 对 $\overline{X'_{i}}$ 和 $\overline{X'_{j}}$ 作类似论证, 可知 \eqref{eq:closures-of-two-quasiprojective-decompositions} 的两个分解确实都是射影簇 $\overline{X}\subseteq\mathbb{P}^{n}$ 的不可约分解.

由射影簇的不可约分解的唯一性 (Proposition/Definition~\ref{prop:projective-irreducible-decomposition}), 有 $r=s$, 并且在必要时重新排序后, 对所有 $i$ 有 $\overline{X_{i}}=\overline{X'_{i}}$. 再次利用 $X_{i}=\overline{X_{i}}\cap X$ 和 $X'_{i}=\overline{X'_{i}}\cap X$ (Exercise~\ref{ex:closure-in-subspace-is-restricted-ambient-closure}), 可知对所有 $i$ 都有 $X_{i}=X'_{i}$. 因而拟射影簇的不可约分解唯一.
\end{Proof}

\BookExerciseHeading
\begin{BookExercises}
  \item\label{ex:closed-subset-equals-topological-closure} 设 $X$ 为拓扑空间, $X\subseteq Y$ 为子集. 证明 $X$ 是闭的, 当且仅当 $X=\overline{X}$.
  \item\label{ex:density-characterized-by-open-sets} 设 $Y$ 为拓扑空间, $X\subseteq Y$ 为子集. 证明 $X$ 在 $Y$ 中稠密, 当且仅当 $Y$ 的每个开集至少包含 $X$ 的一个元素.
  \item\label{ex:closure-commutes-with-finite-unions}
  \begin{enumerate}
    \item 设 $X$ 为拓扑空间, $Y_{1},\ldots,Y_{k}\subseteq X$ 为子集. 证明
    \[
      \overline{Y_{1}\cup\cdots\cup Y_{k}}
      =\overline{Y_{1}}\cup\cdots\cup\overline{Y_{k}}.
    \]
    \item (a) 的结果能否推广到无穷并? 至少有一个包含关系仍成立吗? 证明你的回答, 或给出反例来说明.
  \end{enumerate}
  \item\label{ex:closure-in-subspace-is-restricted-ambient-closure} 设 $Y$ 为拓扑空间, $X\subseteq Y$ 是带子空间拓扑的子集. 对任意子集 $Z\subseteq X$, 证明 $\mathrm{cl}_{X}(Z)=\mathrm{cl}_{Y}(Z)\cap X$.
  \item\label{ex:quasiprojective-variety-open-in-its-closure} 设 $X\subseteq\mathbb{P}^{n}$ 为拟射影簇. 证明 $X$ 是其射影闭包 $\overline{X}$ 的开子集.
  \item\label{ex:affine-closure-preserves-vanishing-ideal} 设 $X\subseteq\mathbb{A}^{n}$ 为一子集, 闭包为 $\overline{X}=\mathrm{cl}_{\mathbb{A}^{n}}(X)$.
  \begin{enumerate}
    \item 证明 $\mathcal{I}_{\mathbb{A}}(X)=\mathcal{I}_{\mathbb{A}}(\overline{X})$.
    \item 推出若 $f,g\in K[x_{1},\ldots,x_{n}]$, 且 $f|_{X}=g|_{X}$ 作为 $X$ 上的函数相等, 则 $f|_{\overline{X}}=g|_{\overline{X}}$.
  \end{enumerate}
  \item\label{ex:irreducible-components-of-dense-open-subvariety} 设 $X$ 是拟射影簇, 不可约分支为 $X_{1},\ldots,X_{r}$. 若 $Y\subseteq X$ 是稠密开子簇, 证明 $X_{1}\cap Y,\ldots,X_{r}\cap Y$ 为 $Y$ 的不可约分支.
  \item\label{ex:irreducibility-criterion-by-nonempty-open-sets} 证明 $X$ 是不可约拓扑空间, 当且仅当 $X$ 的任意两个非空开子集的交都是 $X$ 的非空开子集.
\end{BookExercises}

\par\medskip
{\footnotesize
\noindent\textbf{开放获取.} 本章按照
\href{http://creativecommons.org/licenses/by-nc/4.0/}{Creative Commons Attribution-NonCommercial 4.0 International License}
的条款授权. 该许可允许以任何媒介或格式进行非商业性使用, 分享, 改编, 分发和复制, 但须对原作者和来源作适当署名, 提供 Creative Commons 许可的链接, 并说明是否作了修改.

\noindent
\begin{minipage}[t]{.17\linewidth}
  \vspace{0pt}\includegraphics[width=.9\linewidth]{figures/source-cc-by-nc.png}
\end{minipage}\hfill
\begin{minipage}[t]{.81\linewidth}
  \vspace{0pt}本章的图片或其他第三方材料均纳入本章的 Creative Commons 许可, 除非材料的署名行另有注明. 如果材料未纳入本章的 Creative Commons 许可, 且您拟进行的使用未获法律规定的许可或超出允许使用的范围, 则须直接向著作权持有人取得许可.
\end{minipage}\par}
